\subsection{\texorpdfstring{映射 \((u, v) \mapsto v \circ u\) 的亚连续性}{映射 (u, v) \textbackslash mapsto v \textbackslash circ u 的亚连续性}}

命题 9. --- 设 R、S、T 是三个 Hausdorff 局部凸空间。假设空间 \(\mathcal{L}(\mathbf{R};\mathbf{S})\) 、\(\mathcal{L}(\mathbf{S};\mathbf{T})\) 、\(\mathcal{L}(\mathbf{R};\mathbf{T})\) 分别赋予简单（相应地，紧、有界）收敛拓扑。则双线性映射 \((u, v) \mapsto v \circ u\) （从 \(\mathcal{L}(\mathrm{R}; \mathrm{S}) \times \mathcal{L}(\mathrm{S}; \mathrm{T})\) 到 \(\mathcal{L}(\mathrm{R}; \mathrm{T})\) 中）是 \((\mathfrak{S}, \mathfrak{T})\) -亚连续的，其中 \(\mathfrak{T}\) 是 \(\mathcal{L}(\mathrm{S}; \mathrm{T})\) 的等度连续子集所成的族，而 \(\mathfrak{S}\) 是 \(\mathcal{L}(\mathrm{R}; \mathrm{S})\) 的有限（相应地，紧、有界）子集所成的族。

我们先证明 \((u, v) \mapsto v \circ u\) 是 \(\mathfrak{T}\) -亚连续的。设 H 是 \(\mathcal{L}(S; T)\) 中的一个等度连续集，W 是 T 中 0 的一个邻域，M 是 R 的一个有限（相应地，紧、有界）子集。我们必须证明，存在 S 中 0 的一个邻域 V，使得若 \(u(M) \subset V\) 且 \(v \in H\) ，则 \(v(u(M)) \subset W\) 。但为此，只需 \(v(V) \subset W\) 对一切 \(v \in H\) 成立，而这样的邻域的存在性由 H 的等度连续性得出。

为了看出 \((u, v) \mapsto v \circ u\) 是 \(\mathfrak{S}\) -亚连续的，我们将证明：对 0 的每个邻域 \(W\) （在 \(T\) 中）、每个有限（相应地，紧、有界）子集 \(M\) （在 \(R\) 中）以及每个有限（相应地，紧、有界）子集 \(L\) （在 \(\mathcal{L}(R; S)\) 中），存在一个有限（相应地，紧、有界）子集 \(N\) （在 \(S\) 中） ，使得关系式 \(v(N) \subset W\) 与 \(u \in L\) 蕴含 \(v(u(M)) \subset W\) 。显然，只需证明可取 \(N = \bigcup_{u \in L} u(M)\) ，

即证明只要 L 与 M 是有限（相应地，紧、有界）的，集 N 就是有限（相应地，紧、有界）的。若 L 与 M 都是有限的，这立即可得；若 M 在 R 中有界且 L 在 \(\mathcal{L}(\mathrm{R};\mathrm{S})\) 中有界（对有界收敛拓扑，cf.~III, 第 22 页），这也立即可得。最后，我们证明：若 M 是 R 中的紧集，且 L 对紧收敛拓扑是 \(\mathcal{L}(\mathrm{R};\mathrm{S})\) 中的紧集，则 N 是 S 中的紧集。但若 \(u_{M}\) 是 \(u \in L\) 到 M 上的限制，则映射 \(u \mapsto u_{M}\) （从 L 到由 M 到 S 中的全体连续映射所成的空间 \(\mathcal{C}(\mathrm{M};\mathrm{S})\)（赋一致收敛拓扑）中）是连续的；因此 L 在此映射下的像是紧的，而我们的断言由映射 \((w,x) \mapsto w(x)\) （从 \(\mathcal{C}(\mathrm{M};\mathrm{S}) \times \mathrm{M}\) 到 S 中）的连续性得出（GT, X, § 1, No.~6, 命题 9）。

在下面两个推论中，我们同命题 9 一样假设：空间 \(\mathcal{L}(\mathbf{R};\mathbf{S})\) 、\(\mathcal{L}(\mathbf{S};\mathbf{T})\) 、\(\mathcal{L}(\mathbf{R};\mathbf{T})\) 三者都赋予简单收敛拓扑，或三者都赋予紧收敛拓扑，或三者都赋予有界收敛拓扑。

推论 1. --- 对 \(\mathcal{L}(\mathbf{S};\mathbf{T})\) 的每个等度连续子集 H，映射 \((u,v)\mapsto v\circ u\) （从 \(\mathcal{L}(\mathbf{R};\mathbf{S})\times \mathbf{H}\) 到 \(\mathcal{L}(\mathbf{R};\mathbf{T})\) 中）是连续的。

这由命题 9（III, 第 32 页）与命题 4（III, 第 31 页）立即得出。

推论 2. --- 假设 S 是桶型空间。若序列 \((u_n)\) 收敛于 \(u\) （在 \(\mathcal{L}(\mathbf{R};\mathbf{S})\) 中），且序列 \((v_{n})\) 收敛于 \(v\) （在 \(\mathcal{L}(\mathbf{S},\mathbf{T})\) 中），则序列 \((v_{n}\circ u_{n})\) 收敛于 \(v\circ u\) （在 \(\mathcal{L}(\mathbf{R};\mathbf{T})\) 中） 。

事实上，序列 \((v_{n})\) 在 \(\mathcal{L}(\mathrm{S};\mathrm{T})\) 中简单有界，由于 \(\mathrm{S}\) 是桶型空间（III, 第 25 页, 定理 1），它是等度连续的；于是该推论是推论 1 的结论。

